% ============================================================
% CHAPTER 4 — Project 1: Hello World in 0-Code
% ============================================================
\chapter{Project 1: Hello World in 0-Code}\index{Python}

\section*{What You Will Build}
An interactive Python program that:
\begin{itemize}
  \item Greets the user by name
  \item Calculates the age from the birth year
  \item Generates a random ``fortune of the day''
  \item Has working automated tests
\end{itemize}

\textbf{Estimated time}: 15--20 minutes

% ---------------------------------------------------------
\section{Setting Up the Project (ADLC Phases 0--2)}

\begin{praticabox}[title={Hands-on --- Initial setup}]
\begin{enumerate}
  \item Open VS~Code
  \item Create a new folder: \texttt{hello-0code}
  \item Open the folder in VS~Code (\texttt{File $\rightarrow$ Open Folder})
  \item Create a file \texttt{\_CONTEXT.md} in the root with this content:
\end{enumerate}
\end{praticabox}

\begin{lstlisting}[language={},title={\texttt{\_CONTEXT.md} for Hello 0-Code}]
# Progetto: Hello 0-Code

## Scopo
Un programma Python interattivo che saluta l'utente, calcola
la sua eta' e genera una "fortuna del giorno" casuale.

## Tecnologie
- Python 3.11+
- Solo librerie standard (random, datetime)
- Test con pytest

## Struttura
hello-0code/
+-- _CONTEXT.md
+-- main.py          <- Entry point
+-- greeter.py       <- Logica di saluto e calcolo
+-- fortunes.py      <- Lista di fortune e generatore
+-- tests/
    +-- test_greeter.py
    +-- test_fortunes.py

## Convenzioni
- snake_case per tutto
- Type hints su tutte le funzioni
- Docstring per ogni funzione pubblica
- Print formattato con f-strings

## Vincoli
- Solo librerie standard Python. Nessun pip install.
- L'anno di nascita deve essere validato (1900-anno corrente)
- Il nome non puo' essere vuoto
- Gestisci input non validi con messaggi chiari
\end{lstlisting}

\begin{checkpointbox}
You should have a folder with a single file: \texttt{\_CONTEXT.md}. The AI will generate everything else.
\end{checkpointbox}

% ---------------------------------------------------------
\section{The First Prompt (ADLC Phases 3--4)}

\begin{praticabox}[title={Hands-on --- Generating the code}]
\begin{enumerate}
  \item Open Copilot Chat in Agent Mode (\texttt{Ctrl+Shift+I}, select ``Agent'')
  \item Type this prompt:
\end{enumerate}
\begin{lstlisting}[language={}]
Leggi il file _CONTEXT.md e implementa l'intero progetto
seguendo la struttura e i vincoli definiti.

Il programma deve:
1. Chiedere il nome dell'utente
2. Chiedere l'anno di nascita
3. Calcolare l'eta'
4. Mostrare un saluto personalizzato con l'eta'
5. Mostrare una fortuna del giorno casuale (10+ fortune)

Crea anche i test unitari per greeter.py e fortunes.py.
\end{lstlisting}
\begin{enumerate}
  \setcounter{enumi}{2}
  \item \textbf{Watch the AI at work}: Copilot in Agent Mode will create the files one by one and ask you whether to accept the changes.
  \item \textbf{Accept or correct}: If the code is correct, accept it. If you notice something odd, ask for a revision.
\end{enumerate}
\end{praticabox}

% ---------------------------------------------------------
\section{Analyzing the Output}

Let's look at what the AI should have generated. Do not worry if your output is slightly different --- what matters is that it respects the structure of \texttt{\_CONTEXT.md}.

\subsection*{Expected output: \texttt{greeter.py}}

\begin{lstlisting}[language=Python,title={\texttt{greeter.py}}]
from datetime import datetime


def get_user_name() -> str:
    """Chiede il nome all'utente e lo valida."""
    while True:
        name = input("Come ti chiami? ").strip()
        if name:
            return name
        print("Il nome non puo' essere vuoto. Riprova.")


def get_birth_year() -> int:
    """Chiede l'anno di nascita e lo valida."""
    current_year = datetime.now().year
    while True:
        try:
            year = int(input("In che anno sei nato/a? "))
            if 1900 <= year <= current_year:
                return year
            print(f"L'anno deve essere tra 1900 e {current_year}.")
        except ValueError:
            print("Inserisci un numero valido.")


def calculate_age(birth_year: int) -> int:
    """Calcola l'eta' a partire dall'anno di nascita."""
    return datetime.now().year - birth_year


def greet(name: str, age: int) -> str:
    """Genera il messaggio di saluto personalizzato."""
    return f"Ciao {name}! Hai {age} anni."
\end{lstlisting}

\subsection*{What to look for}

Check that the AI followed the \texttt{\_CONTEXT.md}:

\begin{center}
\begin{tabularx}{\textwidth}{X c}
\toprule
\textbf{Context constraint} & \textbf{Met?} \\
\midrule
snake\_case for functions & \checkmark \\
Type hints & \checkmark \\
Docstrings & \checkmark \\
Year validation (1900--current) & \checkmark \\
Name not empty & \checkmark \\
Standard library only & \checkmark \\
\bottomrule
\end{tabularx}
\end{center}

\begin{tipbox}
This verification is your main job as a 0-code developer. You do not need to write the code, but you do need to know how to \textbf{read and validate} that it respects the requirements.
\end{tipbox}

% ---------------------------------------------------------
\section{Running and Testing (ADLC Phase 5)}

\begin{praticabox}[title={Hands-on --- Running the program}]
Open the integrated terminal and run:
\begin{lstlisting}[language=bash]
python main.py
\end{lstlisting}
Interact with the program:
\begin{lstlisting}[language={}]
Come ti chiami? Marco
In che anno sei nato/a? 1990
Ciao Marco! Hai 36 anni.
La tua fortuna del giorno: "Il codice migliore e'
quello che non devi scrivere."
\end{lstlisting}
\end{praticabox}

\begin{praticabox}[title={Hands-on --- Running the tests}]
\begin{lstlisting}[language=bash]
pip install pytest
python -m pytest tests/ -v
\end{lstlisting}
Expected output:
\begin{lstlisting}[language={}]
tests/test_greeter.py::test_calculate_age PASSED
tests/test_greeter.py::test_greet PASSED
tests/test_greeter.py::test_greet_message_format PASSED
tests/test_fortunes.py::test_get_fortune_returns_string PASSED
tests/test_fortunes.py::test_fortune_not_empty PASSED
========================= 5 passed =========================
\end{lstlisting}
\end{praticabox}

\begin{checkpointbox}
If the program works and the tests pass, you have completed your first project in 0-code.
\end{checkpointbox}

% ---------------------------------------------------------
\section{Iterating and Improving}

The real power of 0-code emerges during iteration.

\begin{praticabox}[title={Hands-on --- Adding a feature}]
In Copilot Chat, type:
\begin{lstlisting}[language={}]
Aggiungi una funzionalita' "zodiac": dopo il saluto, mostra
il segno zodiacale dell'utente basato sull'anno di nascita
(zodiaco cinese). Aggiungi la logica in un nuovo file
zodiac.py con i relativi test. Aggiorna main.py.
\end{lstlisting}
The AI will create \texttt{zodiac.py}, the corresponding tests, and update \texttt{main.py}.
\end{praticabox}

\begin{praticabox}[title={Hands-on --- Fixing a bug}]
If the program behaves unexpectedly, do not edit the code by hand. Describe the problem to Copilot:
\begin{lstlisting}[language={}]
Quando inserisco l'anno 2024 come anno di nascita, il
programma dice che ho 2 anni. Dovrebbe calcolare l'eta'
usando anche mese e giorno, non solo l'anno. Correggi.
\end{lstlisting}
\end{praticabox}

% ---------------------------------------------------------
\section{Lessons from the First Project (ADLC Phase 6)}

\begin{enumerate}
  \item \textbf{\texttt{\_CONTEXT.md} drives everything}: file structure, conventions, constraints --- the AI followed them because you defined them in the context.
  \item \textbf{The initial prompt can be generic}: the details live in the context, while the prompt states the goal.
  \item \textbf{Iteration is free}: adding features or fixing bugs is as fast as describing the problem.
  \item \textbf{You review, the AI implements}: your job is to read the code and verify that it respects the requirements.
\end{enumerate}

\begin{praticabox}[title={Hands-on --- Update the context}]
If the AI made recurring mistakes, add them as constraints:
\begin{lstlisting}[language={}]
## Lezioni Apprese
- Calcolare l'eta' usando datetime.date completo
- I test devono usare fixture, non valori hardcoded
\end{lstlisting}
\end{praticabox}

% ---------------------------------------------------------
\section*{Project Summary}

\begin{center}
\begin{tabularx}{\textwidth}{l X c}
\toprule
\textbf{ADLC Phase} & \textbf{What you did} & \textbf{Time} \\
\midrule
Phase 0--2 & Created \texttt{\_CONTEXT.md} & 5~min \\
Phase 3--4 & Generated code with Copilot & 5~min \\
Phase 5 & Tested and verified & 3~min \\
Iteration & Added zodiac and fixed bugs & 5~min \\
Phase 6 & Updated the context & 2~min \\
\midrule
\textbf{Total} & \textbf{Complete application with tests} & \textbf{$\sim$20~min} \\
\bottomrule
\end{tabularx}
\end{center}
